\chapter{Reproducibility}
\label{app:repro}

The code, the explanations and the result tables are in the GitHub repository; the fine-tuned adapters, the
activations and the cached reconstructions are on Hugging Face (links in Section~\ref{sec:availability}).
\texttt{fetch\_artifacts.py} downloads the large files into place.

\begin{table}[h]
\centering
\small
\setlength{\tabcolsep}{4pt}
\begin{tabular}{@{}ll@{}}
\toprule
Result & Command \\
\midrule
Chapter~\ref{ch:fading}: fade curve, probe, names & \texttt{python fade\_curve.py} \\
Chapters~\ref{ch:fading}, \ref{ch:agreement}: positions, agreement, filter & \texttt{python name\_analysis.py} \\
Chapter~\ref{ch:incentives}: edit test & \texttt{python edit\_test.py} \\
Chapter~\ref{ch:agreement}: confidence & \texttt{python confidence.py} \\
Chapter~\ref{ch:finetuning}: training & \texttt{python finetune.py configs/holdback.json} \\
Chapter~\ref{ch:finetuning}: evaluation & \texttt{python evaluate.py <adapter>} \\
Chapter~\ref{ch:finetuning}: paired comparison & \texttt{python compare.py <names>} \\
Every table and figure here & \texttt{python paper/build\_assets.py} \\
\bottomrule
\end{tabular}
\caption{Which script produces which result.}
\end{table}

All but training and evaluation run on a laptop from the saved files. Seeds are fixed throughout.

\section{Hardware}

Three machines were used (Table~\ref{tab:hardware}). The models run in 16-bit precision on all of them.

\begin{table}[h]
\centering
\footnotesize
\setlength{\tabcolsep}{4pt}
\begin{tabular}{@{}p{0.2\linewidth}p{0.27\linewidth}p{0.3\linewidth}p{0.15\linewidth}@{}}
\toprule
Machine & Specification & Used for & Precision \\
\midrule
Laptop & MacBook Air, Apple M5, 10 cores, 24\,GB memory, macOS 26.6 & Extracting activations, probes, every analysis and table & bfloat16 (MPS) \\
Kaggle notebook & 2 $\times$ NVIDIA Tesla T4, 15\,GB each & Pipeline check; generating the 8{,}640 explanations (batch 32, 70 per minute, 123 minutes) & float16 \\
Rented GPU & 1 $\times$ NVIDIA RTX A6000, 48\,GB; 16 virtual CPUs, 62\,GB memory & Reconstruction run; evaluating the released and reconstruction verbalizers & bfloat16 \\
Rented GPUs & 2 $\times$ NVIDIA A40, 48\,GB, one run per card & All other fine-tuning runs and evaluations; token probabilities & bfloat16 \\
\bottomrule
\end{tabular}
\caption{Hardware.}
\label{tab:hardware}
\end{table}

On the 48\,GB cards generation ran at batch 128, about 800 to 900 explanations per minute. Fine-tuning peaked
at 44\,GB of GPU memory, and a three-round run with its evaluation took about 100 minutes.

\section{Software}

\begin{table}[h]
\centering
\footnotesize
\setlength{\tabcolsep}{4pt}
\begin{tabular}{@{}lll@{}}
\toprule
 & Laptop & Rented GPUs \\
\midrule
Python & 3.12.13 & 3.12 \\
PyTorch & 2.15.0 & 2.8.0 with CUDA 12.8 (Ubuntu 24.04) \\
transformers & 4.57.6 & 4.57 \\
peft & 0.21.2 & 0.21.2 \\
scikit-learn, spaCy, NLTK & 1.9.1, 3.8.18, 3.10.4 & same \\
\bottomrule
\end{tabular}
\caption[Software versions]{Software versions. The Kaggle notebook used PyTorch 2.11.0 with the same
transformers range. \texttt{requirements.txt} pins the laptop environment.}
\label{tab:software}
\end{table}

The three models are fixed to these revisions on Hugging Face:
{\small
\begin{itemize}
\item target model, \texttt{Qwen/Qwen2.5-7B-Instruct}, revision \texttt{a09a3545}
\item verbalizer, \texttt{kitft/nla-qwen2.5-7b-L20-av}, revision \texttt{b8846916}
\item reconstructor, \texttt{kitft/nla-qwen2.5-7b-L20-ar}, revision \texttt{e2c9e57e}
\end{itemize}}
Fine-tuning used a LoRA of rank 16
(scale 32, dropout 0.05) on the four attention projections, AdamW with learning rate $10^{-4}$, batch size 8
and one pass over the chosen explanations per round.

\chapter{Layer Indexing}
\label{app:layer}

The NLA for Qwen2.5-7B is trained on ``layer 20'', meaning the output of transformer block 20. In the
Hugging Face convention \texttt{hidden\_states[0]} is the embedding output, so the output of block 20 is
\texttt{hidden\_states[21]}. The usage example published with the models reads \texttt{hidden\_states[20]},
one block early.

Nothing fails when this happens. The explanations still read sensibly, because neighbouring layers are
similar. Only the numbers give it away: on the same passages, reading one block early gave a mean cosine of
0.846 and an activation norm of 95.7, against 0.893 and 112.7 at the right layer, and the reference norm
for this layer is 111.3. I found the error when the first check against a known reconstruction quality came
out low, and all results in this report use the corrected layer.

\chapter{Full Tables}
\label{app:tables}

\begin{table}[h]
\centering
\footnotesize
\setlength{\tabcolsep}{4pt}
\begin{tabular}{@{}llllll@{}}
\toprule
\shortstack[l]{Tokens after\\entity} & \shortstack[l]{Keep a name if\\given by} & \shortstack[l]{Wrong names\\removed} & \shortstack[l]{True names\\kept} & Precision & Names left \\
\midrule
0 & None & 0 & 100 & 34.6 {\scriptsize[32.0, 37.5]} & 1{,}559 \\
0 & At least 2 of 4 & 70.4 {\scriptsize[66.6, 74.0]} & 100.0 {\scriptsize[100.0, 100.0]} & 64.1 {\scriptsize[60.5, 68.1]} & 842 \\
0 & At least 3 of 4 & 84.9 {\scriptsize[82.5, 87.2]} & 100.0 {\scriptsize[100.0, 100.0]} & 77.8 {\scriptsize[74.5, 81.3]} & 694 \\
0 & All 4 & 92.4 {\scriptsize[90.8, 94.0]} & 99.8 {\scriptsize[99.4, 100.0]} & 87.5 {\scriptsize[84.8, 90.3]} & 616 \\
1 & None & 0 & 100 & 19.1 {\scriptsize[16.7, 21.8]} & 1{,}883 \\
1 & At least 2 of 4 & 69.3 {\scriptsize[65.4, 73.0]} & 83.9 {\scriptsize[79.8, 87.6]} & 39.2 {\scriptsize[34.4, 44.3]} & 770 \\
1 & At least 3 of 4 & 82.8 {\scriptsize[80.0, 85.5]} & 74.4 {\scriptsize[69.7, 79.0]} & 50.6 {\scriptsize[45.0, 56.3]} & 530 \\
1 & All 4 & 92.3 {\scriptsize[90.7, 93.9]} & 64.2 {\scriptsize[57.9, 70.1]} & 66.4 {\scriptsize[59.7, 73.0]} & 348 \\
4 & None & 0 & 100 & 18.1 {\scriptsize[15.0, 21.7]} & 975 \\
4 & At least 2 of 4 & 77.0 {\scriptsize[72.7, 81.1]} & 70.5 {\scriptsize[64.2, 76.5]} & 40.3 {\scriptsize[33.5, 47.7]} & 308 \\
4 & At least 3 of 4 & 88.5 {\scriptsize[85.7, 91.1]} & 57.4 {\scriptsize[52.1, 63.0]} & 52.3 {\scriptsize[44.7, 60.5]} & 193 \\
4 & All 4 & 94.9 {\scriptsize[93.2, 96.4]} & 38.6 {\scriptsize[32.7, 44.4]} & 62.4 {\scriptsize[53.0, 71.7]} & 109 \\
16 & None & 0 & 100 & 11.7 {\scriptsize[8.7, 15.4]} & 631 \\
16 & At least 2 of 4 & 80.6 {\scriptsize[75.9, 84.8]} & 60.8 {\scriptsize[49.4, 71.6]} & 29.4 {\scriptsize[21.3, 38.1]} & 153 \\
16 & At least 3 of 4 & 92.6 {\scriptsize[89.5, 95.3]} & 39.2 {\scriptsize[27.2, 50.6]} & 41.4 {\scriptsize[27.9, 55.4]} & 70 \\
16 & All 4 & 97.1 {\scriptsize[95.6, 98.4]} & 17.6 {\scriptsize[10.1, 24.7]} & 44.8 {\scriptsize[26.9, 62.1]} & 29 \\
1, 4 and 16 & None & 0 & 100 & 17.5 {\scriptsize[15.1, 20.1]} & 3{,}489 \\
1, 4 and 16 & At least 2 of 4 & 73.6 {\scriptsize[69.9, 77.3]} & 77.2 {\scriptsize[73.5, 80.6]} & 38.3 {\scriptsize[33.5, 43.4]} & 1{,}231 \\
1, 4 and 16 & At least 3 of 4 & 86.3 {\scriptsize[83.9, 88.6]} & 65.2 {\scriptsize[61.6, 68.9]} & 50.2 {\scriptsize[44.6, 56.2]} & 793 \\
1, 4 and 16 & All 4 & 94.0 {\scriptsize[92.7, 95.2]} & 51.1 {\scriptsize[46.3, 55.9]} & 64.2 {\scriptsize[57.7, 70.5]} & 486 \\
\bottomrule
\end{tabular}

\caption[The agreement filter at each distance]{The agreement filter at each distance (\%, 95\% intervals;
natural text). The pooled rows are Table~\ref{tab:filter}.}
\end{table}

\begin{table}[h]
\centering
\footnotesize
\setlength{\tabcolsep}{4pt}
\begin{tabular}{@{}ccllcl@{}}
\toprule
\shortstack[l]{Tokens after\\entity} & \shortstack[l]{Wrong names,\\released} & Change & Relative & \shortstack[l]{True-name\\rate, released} & \shortstack[l]{Change\\(points)} \\
\midrule
0 & 0.719 & $-$0.051 {\scriptsize[$-$0.088, $-$0.014]} & $-$7.2\% & 100.0\% & $+$0.0 {\scriptsize[$+$0.0, $+$0.1]} \\
1 & 1.097 & $-$0.091 {\scriptsize[$-$0.130, $-$0.054]} & $-$8.3\% & 53.8\% & $-$0.9 {\scriptsize[$-$2.6, $+$0.6]} \\
4 & 0.517 & $-$0.055 {\scriptsize[$-$0.088, $-$0.022]} & $-$10.7\% & 21.7\% & $-$0.6 {\scriptsize[$-$2.1, $+$1.0]} \\
16 & 0.334 & $-$0.025 {\scriptsize[$-$0.062, $+$0.006]} & $-$7.6\% & 7.5\% & $+$0.6 {\scriptsize[$-$0.5, $+$1.9]} \\
\bottomrule
\end{tabular}

\caption[Hold back at each distance]{Hold back (first run) against the released verbalizer at each distance. Natural
text, held-out passages; 95\% intervals from a paired bootstrap over entities.}
\end{table}

\chapter{Examples}
\label{app:examples}

Four more explanations from the released verbalizer, with the relevant part of each passage. True names are
in bold; wrong names are bold and underlined.

\begin{sample}{One token after \emph{eagle}: a wrong animal}
Passage: \ldots Superimposed on the cross is a silver crowned \true{eagle} for all grades.\\[2pt]
Explanation: \ldots Final token ``for'' ends an incomplete clause (``a black \wrong{lion} for''), part of a
listing of insignia details (``A gold \wrong{lion} crest with a crown for'') \ldots
\end{sample}

\begin{sample}{One token after \emph{bread}: a wrong food}
Passage: \ldots Guilmon possesses a taste for \true{bread}, often being distracted by the thought of food.\\[2pt]
Explanation: \ldots ``Momo has a strong craving for food, particularly sweets and ice cream \ldots\ He also has
a fondness for \wrong{carrots}'' \ldots\ Final token ``\wrong{sugar},'' ends an incomplete clause \ldots
\end{sample}

\begin{sample}{Four tokens after \emph{wolf}: no name, which is the faithful answer}
Passage: \ldots the seahorse and the ship refer to Belfast's significant maritime history. The \true{wolf} may
be a tribute to the city's founder \ldots\\[2pt]
Explanation: \ldots ``\ldots may include a reference to the state of New York or a tribute''
suggests a third possible meaning or origin story \ldots
\end{sample}

\begin{sample}{Sixteen tokens after \emph{Shakespeare}: the name survives}
Passage: \ldots a line by William \true{Shakespeare} penned in 1594 in Love's Labour's Lost, Act IV\\[2pt]
Explanation: \ldots ``As a fact the \true{Shakespeare}'s play `Titanium' refers to his famous passage in
Act IV'' references a specific quote \ldots
\end{sample}
